\documentclass[11pt]{article}
\usepackage{amsmath,amsthm,amssymb}
\usepackage[margin=1in]{geometry}

\newtheorem{theorem}{Theorem}
\newtheorem{lemma}[theorem]{Lemma}
\newtheorem{proposition}[theorem]{Proposition}
\newtheorem{corollary}[theorem]{Corollary}
\theoremstyle{definition}
\newtheorem{definition}[theorem]{Definition}
\theoremstyle{remark}
\newtheorem{remark}[theorem]{Remark}

\newcommand{\rank}{\operatorname{rank}}
\newcommand{\im}{\operatorname{im}}

\title{The Contraction Lemma and the Per-Irrep Upper Bound\\
for the Staircase Product}
\author{Clio}
\date{April 15, 2026}

\begin{document}
\maketitle

\begin{abstract}
We prove that for the staircase symmetrizing product
$\Pi = \prod (1+s_{i_j})/2$ of the longest element $w_0 \in S_n$,
the image of $\Pi$ on any irreducible representation $V_\lambda$
satisfies $\rank(\Pi|_{V_\lambda}) \leq \dim(V_\lambda^H)$, where
$H = \langle s_1, s_3, s_5, \ldots \rangle$.
The key tool is a contraction lemma showing that adjacent projections
cannot create pure negative eigenvectors.
Combined with the total rank identity
$\sum_\lambda \dim(V_\lambda) \cdot \rank(\Pi|_{V_\lambda})
= n!/2^{\lfloor n/2 \rfloor}$
(verified computationally for $n \leq 9$),
this gives $\rank(\Pi|_{V_\lambda}) = \dim(V_\lambda^H)$ for all $\lambda$.
\end{abstract}

\section{Setup}

We work in Young's orthogonal representation throughout.
Each simple transposition $s_i$ acts as an involution $S_i$ with $S_i^2 = I$
and $S_i = S_i^T$ (symmetric). The projection $P_i = (I + S_i)/2$ is
an \textbf{orthogonal} projection onto $E_i^+ = \ker(S_i - I)$.

The staircase product is $\Pi = \prod_{j=1}^{N} P_{i_j}$ where the
staircase word is $(1, \; 2,1, \; 3,2,1, \; \ldots, \; n{-}1, \ldots, 1)$,
organized into blocks: block $k$ applies $P_k, P_{k-1}, \ldots, P_1$.

Let $H = \langle s_1, s_3, s_5, \ldots \rangle \leq S_n$ and
$e_H = \prod_{j \text{ odd}} P_j$ (the orthogonal projection onto $V^H$,
since the odd generators pairwise commute).

\section{The Contraction Lemma}

\begin{lemma}[Adjacent Disjointness]\label{lem:adj}
For adjacent transpositions $s_i, s_{i+1}$ in any representation of $S_n$:
$E_i^+ \cap E_{i+1}^- = \{0\}$ and $E_i^- \cap E_{i+1}^+ = \{0\}$.
\end{lemma}

\begin{proof}
Standard: decompose into $S_3 = \langle s_i, s_{i+1} \rangle$ irreps.
In the trivial irrep both eigenvalues are $+1$; in the sign irrep both are $-1$;
in the standard 2D irrep, $E_i^+$ and $E_{i+1}^-$ are distinct 1D subspaces.
\end{proof}

\begin{lemma}[Contraction Lemma]\label{lem:contraction}
For any representation of $S_n$ (in orthogonal form), for any $1 \leq j \leq n-2$:
\[
P_{j-1}(E_j^+) \cap E_j^- = \{0\},
\]
where we interpret $P_0 = P_1$ if $j=1$ (vacuously true since $P_1(E_1^+) = E_1^+ \cap E_1^- = \{0\}$).

More precisely: if $v \in E_j^+$ and $w = P_{j-1}(v)$, then $w \notin E_j^-$
(unless $v = 0$).
\end{lemma}

\begin{proof}
Suppose $v \in E_j^+$ with $v \neq 0$, and $w = P_{j-1}(v) \in E_j^-$.
Then $w = \frac{1}{2}(v + s_{j-1} v)$ and $s_j w = -w$.

Expanding the condition $s_j w = -w$:
\[
\frac{1}{2}(s_j v + s_j s_{j-1} v) = -\frac{1}{2}(v + s_{j-1} v).
\]
Since $v \in E_j^+$: $s_j v = v$. So:
\[
\frac{1}{2}(v + s_j s_{j-1} v) = -\frac{1}{2}(v + s_{j-1} v).
\]
Rearranging:
\[
2v + s_{j-1} v + s_j s_{j-1} v = 0,
\quad\text{i.e.,}\quad
v = -\frac{1+s_j}{2} s_{j-1} v = -P_j(s_{j-1} v).
\]

Now apply the \textbf{contraction inequality}. Since $P_j$ is an
orthogonal projection and $s_{j-1}$ is an isometry:
\[
\|v\| = \|P_j(s_{j-1} v)\| \leq \|s_{j-1} v\| = \|v\|.
\]
All inequalities are therefore equalities. The equality
$\|P_j(s_{j-1} v)\| = \|s_{j-1} v\|$ holds if and only if
$s_{j-1} v \in \im(P_j) = E_j^+$.

With $s_{j-1} v \in E_j^+$: $P_j(s_{j-1} v) = s_{j-1} v$, so
$v = -s_{j-1} v$, which gives $v \in E_{j-1}^-$.

But $v \in E_j^+ \cap E_{j-1}^- = \{0\}$ by adjacent disjointness
(Lemma~\ref{lem:adj}). This contradicts $v \neq 0$.
\end{proof}

\section{Image Avoidance}

\begin{proposition}[Image Avoidance]\label{prop:avoidance}
For the staircase product $\Pi$ acting on any representation $V_\lambda$
of $S_n$:
\[
\im(\Pi) \cap E_j^- = \{0\} \quad \text{for all } j = 1, \ldots, n-2.
\]
\end{proposition}

\begin{proof}
Fix $j \in \{1, \ldots, n-2\}$. The last block of the staircase
(block $n-1$) applies $P_{n-1}, P_{n-2}, \ldots, P_1$ in that order.
Since $j \leq n-2$, both $P_j$ and $P_{j-1}$ appear in this block
(with $P_j$ applied before $P_{j-1}$).

We track the $E_j^-$ content of the image through the last block.

\textbf{Step 1: After $P_j$.}
Since $P_j = (I + S_j)/2$ is the orthogonal projection onto $E_j^+$,
the image after $P_j$ lies in $E_j^+$.

\textbf{Step 2: After $P_{j-1}$.}
By the Contraction Lemma (Lemma~\ref{lem:contraction}):
$P_{j-1}(E_j^+) \cap E_j^- = \{0\}$.
So after $P_{j-1}$, the image contains no vector in $E_j^-$.

\textbf{Step 3: After $P_{j-2}, \ldots, P_1$.}
For each $i \leq j-2$: $|i - j| \geq 2$, so $s_i$ and $s_j$ commute.
Therefore $P_i$ commutes with $S_j$ and preserves the eigenspaces
$E_j^+$ and $E_j^-$. In particular, $P_i$ cannot create new vectors
in $E_j^-$ from vectors not in $E_j^-$.

Combining Steps 1--3: the image at the end of the last block (which is
$\im(\Pi)$) has no vector in $E_j^-$.
\end{proof}

\section{The Per-Irrep Upper Bound}

\begin{theorem}[Per-Irrep Upper Bound]\label{thm:upper}
For every partition $\lambda \vdash n$:
\[
\rank(\Pi|_{V_\lambda}) \leq \dim(V_\lambda^H).
\]
\end{theorem}

\begin{proof}
By Proposition~\ref{prop:avoidance}, $\im(\Pi) \cap E_j^- = \{0\}$
for all $j$. In particular, for each odd $j \in \{1, 3, 5, \ldots\}$
(the generators of $H$):
\[
\im(\Pi|_{V_\lambda}) \cap E_j^- = \{0\}.
\]
This means $P_j$ is injective on $\im(\Pi|_{V_\lambda})$
(since $\ker(P_j) = E_j^-$). Therefore:
\[
\rank(\Pi \cdot P_j |_{V_\lambda}) = \rank(\Pi|_{V_\lambda})
\quad \text{for each odd } j.
\]

Since the odd generators pairwise commute ($|j - j'| \geq 2$ for
distinct odd $j, j'$), the projections $P_j$ for odd $j$ pairwise commute
and their product is $e_H = \prod_{j \text{ odd}} P_j$, the projection
onto $V_\lambda^H$.

Applying each $P_j$ successively (each injective on the current image):
\[
\rank(\Pi|_{V_\lambda})
= \rank(\Pi \cdot P_1 |_{V_\lambda})
= \rank(\Pi \cdot P_1 \cdot P_3 |_{V_\lambda})
= \cdots
= \rank(\Pi \cdot e_H |_{V_\lambda}).
\]

Finally, by the rank inequality for compositions:
\[
\rank(\Pi \cdot e_H |_{V_\lambda})
\leq \rank(e_H |_{V_\lambda})
= \dim(V_\lambda^H). \qedhere
\]
\end{proof}

\section{The H-Invariant Theorem}

\begin{theorem}[H-Invariant Theorem]\label{thm:main}
Assume the total rank identity: in the regular representation of $S_n$,
\[
\rank(\Pi) = \frac{n!}{2^{\lfloor n/2 \rfloor}}.
\]
Then for every partition $\lambda \vdash n$:
\[
\rank(\Pi|_{V_\lambda}) = \dim(V_\lambda^H).
\]
\end{theorem}

\begin{proof}
By Frobenius reciprocity (or direct double coset counting):
\[
\sum_{\lambda \vdash n} \dim(V_\lambda) \cdot \dim(V_\lambda^H)
= \frac{n!}{|H|}
= \frac{n!}{2^{\lfloor n/2 \rfloor}}.
\]

By the total rank identity, the regular representation rank decomposes as:
\[
\sum_{\lambda \vdash n} \dim(V_\lambda) \cdot \rank(\Pi|_{V_\lambda})
= \frac{n!}{2^{\lfloor n/2 \rfloor}}
= \sum_{\lambda \vdash n} \dim(V_\lambda) \cdot \dim(V_\lambda^H).
\]

By Theorem~\ref{thm:upper}, each term satisfies
$\rank(\Pi|_{V_\lambda}) \leq \dim(V_\lambda^H)$.
Since the sums are equal and every term satisfies $\leq$, every
term must satisfy equality:
\[
\rank(\Pi|_{V_\lambda}) = \dim(V_\lambda^H)
\quad \text{for all } \lambda \vdash n. \qedhere
\]
\end{proof}

\begin{remark}[Status of the total rank identity]
The total rank identity $\rank(\Pi) = n!/2^{\lfloor n/2 \rfloor}$
is verified computationally for $n \leq 9$ (all irreps, floating-point
with tolerance $10^{-10}$). An algebraic proof would complete the
H-Invariant Theorem unconditionally. The image basis conjecture
($\im(\Pi) = \operatorname{span}\{e_w : w \text{ descent-paired}\}$
in the regular representation) implies this identity and is verified
for $n \leq 5$.
\end{remark}

\section{Computational Verification}

\begin{enumerate}
\item \textbf{Contraction Lemma} (Lemma~\ref{lem:contraction}):
Verified directly for all irreps of $S_n$, $n = 3, \ldots, 9$.
In every case, $P_{j-1}(E_j^+) \cap E_j^- = \{0\}$.

\item \textbf{Image Avoidance} (Proposition~\ref{prop:avoidance}):
Verified for all irreps of $S_n$, $n = 3, \ldots, 9$.
$\im(\Pi) \cap E_j^- = \{0\}$ for \emph{all} $j$ (not just odd $j$).

\item \textbf{Per-irrep equality}: $\rank(\Pi|_{V_\lambda}) = \dim(V_\lambda^H)$
verified for all $\lambda \vdash n$, $n \leq 9$ (exact arithmetic) and
$n \leq 16$ (floating-point).

\item \textbf{S$_5$ Lemma} (bonus): $\ker(P_a P_c) \cap E_b^+ \cap E_d^+ = \{0\}$
for $S_5 = \langle s_1, s_2, s_3, s_4 \rangle$ with $a = s_2, b = s_3,
c = s_4, d = s_1$. Verified in all 7 irreps of $S_5$.
\end{enumerate}

\end{document}
